\section{A common family before a polynomial closing parameter}
\label{sec:common_closing_polynomial_family}

Retain the two windows, the native middle core, the interval $J$, and
$D=(qk)^{18\ell+2\lfloor k/\ell\rfloor}$ of
Lemma~\ref{thm:physical_polynomial_tensor_family}.
Fix $m\geq0$ and closing matrices $C_0,\ldots,C_m\in\C^{B\times B}$.
The site tensors are fixed; only their closing matrix varies:
\begin{align}
    C(z)&=\sum_{a=0}^m z^a C_a,\notag\\
    Z(z)_{\sigma,\tau}
      &=\operatorname{tr}\left(C(z)
          \prod_{u=0}^{|K|-1}A_u(\sigma_u,\tau_u)\right).
    \label{eq:polynomial_closing_matrix}
\end{align}
Let $P_M(z)$ be the identity extension of $Z(z)$ to the original chain.

\begin{lemma}[One family for all closing parameters]
    \label{thm:common_closing_polynomial_family}
    \lean{TNLean.MPS.AreaLaw.exists_common_physical_polynomial_closing_family_fin}
    \leanok
    Under the geometric and support hypotheses
    (\ref{eq:physical_tensor_hypotheses}), there are a nonnegative integer
    $N$ and matrices $W_0,\ldots,W_{N-1}$ on the retained interval such that
    \begin{align}
        N&\leq(mk+1)D^2.
        \label{eq:common_closing_family_count}
    \end{align}
    Each $W_i$, in the paired physical alphabet and increasing site order,
    has an open-boundary matrix product representation with all bonds at
    most $(\max\{1,B\})^{2k}$.

    This one family is chosen before $z\in\C$, before the supported
    exterior operators $P_L,P_R$, and before the polynomial $P$ of degree
    at most $k$. For
    \begin{align}
        H_z=P_L+P_M(z)+P_R+
          \sum_{r=0}^{2\ell-1}(W^L_r+W^R_r),
        \label{eq:common_closing_hamiltonian}
    \end{align}
    there are complementary matrices $M_i$ such that
    \begin{align}
        P(H_z)&=\left(\sum_{i=0}^{N-1}M_i\otimes W_i\right)^{\mathrm{phys}},
        \label{eq:common_closing_matrix_decomposition}\\
        \mathcal U P(H_z)\mathcal U^*
          &=\sum_{i=0}^{N-1}W_i\otimes M_i.
        \label{eq:common_closing_conjugation}
    \end{align}
    The same canonical coordinate isometry $\mathcal U$ is used for every
    parameter and every polynomial. No Hermiticity assumption is imposed
    on the closing matrices or on the interpolation nodes.
\end{lemma}
\begin{proof}
    \leanok
    \uses{thm:physical_aeval_fin_middle_family,thm:physical_polynomial_tensor_family}
    For each of the $mk+1$ fixed integer nodes, apply the common physical
    family construction to the corresponding closing matrix. Take the
    disjoint union of these families. The count adds, whereas each member
    retains its original matrix product representation.

    Hold $P_L,P_R$ fixed. Linearity in the closing matrix makes every
    entry of $H_z$ a polynomial in $z$ of degree at most $m$, independently
    of the number of sites. Every entry of $P(H_z)$ therefore has degree
    at most $mk$, even though the coefficient matrices need not commute.
    Lagrange interpolation expresses $P(H_z)$ as a scalar linear
    combination of its values at the fixed nodes. Absorb these scalar
    weights into the complementary factors $M_i$. This gives
    (\ref{eq:common_closing_matrix_decomposition}) with the unchanged
    middle family. The canonical tensor-coordinate identity gives
    (\ref{eq:common_closing_conjugation}).
\end{proof}

The exterior operators are held fixed during the interpolation. They may
subsequently be chosen as arbitrary functions of the requested physical
parameter; in particular, a separate scalar energy shift can remain in
$P_L$. The result is an algebraic ingredient for
\cite[Proposition~5.7]{Arad2016RigorousRG}, not a spectral or running-time
estimate.
